\section{Primes with a controlled predecessor in a fixed progression}
\label{sec:prime-construction}
\label{sec:construction}

We prove Proposition~\ref{prop:controlled}.  Fix its data \(M,c,u\), so
\begin{equation}\label{con:progression}
 M\equiv0\pmod8,\qquad c\in\{2,4\},\qquad
 (u,M)=1,\qquad c\mid u-1,\qquad
 \left(\frac{u-1}{c},\frac Mc\right)=1.
\end{equation}
We first assign nonnegative weights to integers \(d=crQ+1\) in the
specified progression.  A lower sieve bound supplies integers with
no small prime factor.  The marked Type~II estimate then allows us
to bound the mass of composites according to their least prime
factor.  Factors close to the square root require a separate upper
sieve bound.  Keeping the constant in that last bound independent
of the number of marks will permit the final parameter choices.

\subsection{The weights and the Type~II estimate}

Use the parameters $L,W,V$, the prime groups $\mathcal P_i$, and
the marks $\mathcal W$ of
\eqref{eq:sieve-basic-parameters}, \eqref{eq:marked-prime-groups},
and \eqref{eq:marked-weight}, with $q=1/2$.
All prime-indexed sums and products below are over rational primes.
The integer $K\ge1$ and exponents
$0.1<a_1<\cdots<a_K<0.2$ are fixed before $x$ grows; their eventual
choice will follow the other sieve parameters.
The groups are disjoint for large \(x\), and Mertens' theorem gives
\(V_i=\log2+o(1)\).  Their primes are then coprime to \(M\) and lie
in \((\exp(L^{0.1}),\exp(L^{0.3}))\).

The factor \(\omega_i(h)\) counts choices of a prime divisor of
\(h\) from the \(i\)th group. Thus \(\mathcal W(h)\) vanishes
unless every group contributes, while
\(\omega_i(h)2^{1-\omega_i(h)}\) remains bounded when several
primes from one group divide \(h\). These choices are the marks
in the Type~II estimate below.

Call a positive integer a \emph{group integer} if all its prime
divisors belong to \(\bigcup_i\mathcal P_i\), and include \(1\).
Use the complete group-prime part $h_{\mathcal P}$ defined in
\eqref{eq:complete-band-part}. Fix a smooth function $\Psi$
compactly supported in \((1,2)\), with \(0\le\Psi\le1\) and
\(\int\Psi>0\), and define
\begin{equation}\label{con:weight}
 \begin{split}
 F(h)&=\one_{\{h/h_{\mathcal P}=cQ\text{ for a prime }Q>x^{0.9}\}},\\
 w(d)&=\one_{\{d\equiv u\pmod M\}}\Psi(d/x)
                                  F(d-1)\mathcal W(d-1)\qquad(d\ge2).
 \end{split}
\end{equation}
Set \(w(1)=0\).  Since \((ph)_{\mathcal P}=p h_{\mathcal P}\) for
every group prime \(p\), including when \(p\mid h\), we have
\begin{equation}\label{con:invariance}
 F(ph)=F(h)\qquad(h\ge1,\ p\in\textstyle\bigcup_i\mathcal P_i).
\end{equation}
On the support of \(w\), the representation \(d=crQ+1\) is unique:
\(r=(d-1)_{\mathcal P}\), since neither \(c\) nor \(Q\) has a
group prime factor.  In particular,
\begin{equation}\label{con:support}
 r\le x^{0.11},\qquad \mathcal W(d-1)=\mathcal W(r),\qquad
 0\le w(d)\le B_K
\end{equation}
for a constant \(B_K\).  The last assertion follows from
\(j2^{1-j}\le1\) for each integer \(j\ge1\); in fact
\(0\le\mathcal W(h)\le B_K\) for every \(h\ge1\).

Write \(X_0=\sum_d w(d)\) for the total mass.  Our aim is to prove
that the prime mass is \(\gg X_0/L\), while \(X_0\gg x/L\).
The bound \(w\le B_K\) will then give the required count of primes.
The marked estimate supplies the comparison needed to subtract
composite mass.

Use the rough-number density $D_\gamma(s)$ from
\eqref{eq:rough-density-definition}.  For fixed \(\gamma>0\) and integers
\(m>x^\gamma\), put
\begin{equation}\label{con:proxies}
 R_\gamma(m)=\one_{\{P^-(m)>x^\gamma\}},\qquad
 B_\gamma(m)=\frac{D_\gamma(\log m/L)}{LV(W)}
                         \one_{\{P^-(m)>W\}}.
\end{equation}
We will replace \(R_\gamma\) by \(B_\gamma\) inside sums weighted
by \(w(mn)\): the factor \(D_\gamma/(LV(W))\) compensates for
retaining only \(W\)-roughness. Lemma~\ref{lem:marked-coefficient-transfer}
and the character-transform bound proved in
Proposition~\ref{prop:rough-type-ii} permit this replacement after
the fixed congruence is expanded in characters. We verify that
application next. All Dirichlet characters retain the zero extension
off the units used in Section~\ref{sec:type-ii}.

\begin{lemma}[Type~II estimate in the progression]\label{con:type-ii}
Fix $\delta>0$ and $0<\gamma<s_-<s_+<1$. Let
$H_m,H_n\ge x^\delta$, with $X=H_mH_n\asymp x$ and fixed
comparison constants, and let $I\subseteq[H_m,2H_m]$ be an interval.
Suppose
\begin{equation}\label{con:exponent-range}
 [\log H_m/L,\log(2H_m)/L]\subseteq[s_-,s_+].
\end{equation}
Let \(|b_n|\le1\) be supported on \(W\)-rough integers in
\([H_n,2H_n]\).  For sufficiently large \(K\),
\begin{equation}\label{con:type-ii-bound}
 \left|\sum_{m\in I,n}
       \bigl(R_\gamma(m)-B_\gamma(m)\bigr)b_n w(mn)\right|
       \ll XL^{-6}.
\end{equation}
This is uniform in the dyads, interval, and coefficients.  The
lower bound on \(K\) is independent of the \(a_i\). Constants and
thresholds may depend on all fixed parameters, including the band
exponents.
\end{lemma}

\begin{proof}
The congruence in \(w\) cannot be incorporated into \(F\) while
preserving \eqref{con:invariance}.  Instead, character orthogonality
and \((u,M)=1\) give
\[
 \one_{\{mn\equiv u\pmod M\}}
 =\frac1{\varphi(M)}\sum_{\lambda\bmod M}
             \overline{\lambda(u)}\lambda(m)\lambda(n).
\]
By the proof of Proposition~\ref{prop:rough-type-ii}, the sequence
$\alpha_m=\one_I(m)m^{iv}(R_\gamma(m)-B_\gamma(m))$ satisfies
\eqref{eq:type-ii-long-coefficient-interface} for every fixed
$A_0,B,A,C>0$, every character of modulus at most $L^{A_0}$,
and every $|t|\le2XL^B$, uniformly for $|v|\le L^C$ and every
interval $I$ under the present fixed exponent hypotheses.
If $\alpha$ satisfies that bound, so does
\(\alpha_m\lambda(m)\): the product \(\lambda\chi\), with its
zero extension, is a character modulo \(\operatorname{lcm}(M,k)\),
which is at most \(ML^{A_0}\le L^{A_0+1}\) for large \(x\).

To separate the smooth factor, write
\[
 \widehat\psi(v)=\frac1{2\pi}\int_{\R}\Psi(\e^s)\e^{-ivs}\,ds.
\]
Fourier inversion yields
\[
 \Psi(mn/x)=\int_{\R}\widehat\psi(v)x^{-iv}m^{iv}n^{iv}\,dv.
\]
For \(|v|\le L\), apply Lemma~\ref{lem:marked-coefficient-transfer}
to the invariant function $F$ in \eqref{con:weight}, with
\(C=2\), \(D_*=6\), and
\[
 \alpha_m=m^{iv}\lambda(m)
                   (R_\gamma(m)-B_\gamma(m)),\qquad
 \beta_n=b_n\lambda(n)n^{iv}.
\]
Here the formula for \(\alpha_m\) applies on \(I\), and
\(\alpha_m=0\) elsewhere.
Both sequences are \(W\)-rough and bounded by \(L^2\) for large
\(x\): \(x^\gamma>W\), \(D_\gamma\) is bounded on the fixed
exponent interval, and \(LV(W)\asymp L^{0.76}\).  The finite
number of characters is fixed with \(M\).  The integral of
\(|\widehat\psi|\) is finite, whereas its tail beyond \(L\) is
\(O_A(L^{-A})\) for every fixed \(A\).  On that tail the bilinear
sum is \(O_K(X)\) by \eqref{con:support}.  Integration proves
\eqref{con:type-ii-bound}.
\end{proof}

All parameters in the remainder of the section are fixed before
\(x\) tends to infinity.  Unless an independence is explicitly
asserted, constants, thresholds, and rates of convergence may depend
on those fixed parameters.  The estimates are valid for arbitrary
fixed marks, subject to the lower bounds on \(K\) in the Type~II
applications.  We will choose the parameters together at the end.

\subsection{Mass and distribution in progressions}

Let \(r\) range over group integers, and define
\begin{equation}\label{con:mass-definitions}
 \begin{split}
 J_0&=\sum_r\frac{\mathcal W(r)}r,\\
 A_r&=\sum_{\substack{Q>x^{0.9}\text{ prime}\\crQ+1\equiv u\pmod M}}
                         \Psi((crQ+1)/x),\qquad
 X_0=\sum_d w(d)=\sum_r\mathcal W(r)A_r.
 \end{split}
\end{equation}
We shall show that \(X_0\) is of order \(xJ_0/L\), and that the
family at each small \(r\) has enough distribution to apply the
block sieve.  The harmonic sum makes this possible without constants
depending on the pointwise bound \(B_K\).

\begin{lemma}[Harmonic mass]\label{con:harmonic}
For each group put
\[
 E_i(t)=\prod_{p\in\mathcal P_i}\left(1+\frac{t}{p-1}\right).
\]
The harmonic mass has the exact Euler-product expression
\begin{equation}\label{con:harmonic-identity}
 \begin{aligned}
 J_0&=\prod_{i=1}^K\frac{E_i'(1/2)}{V_i}\\
    &=\prod_{i=1}^K\left\{
       \frac{E_i(1/2)}{V_i}
       \sum_{p\in\mathcal P_i}\frac1{p-1+1/2}\right\}.
 \end{aligned}
\end{equation}
For every fixed \(\eps>0\),
\begin{equation}\label{con:harmonic-bounds}
 1\le J_0\ll_K1,\qquad
 \sum_{r>x^\eps}\frac{\mathcal W(r)}r
       \ll_K\exp(-\eps L^{1-a_K}).
\end{equation}
Consequently the part of \(X_0\) with \(r>x^\eps\) is
\(O_K(x\exp(-\eps L^{1-a_K}))\).
\end{lemma}

\begin{proof}
At a prime \(p\), summing over its positive exponents contributes
\(\sum_{j\ge1}p^{-j}=1/(p-1)\). Consequently \(E_i(t)\) is the
harmonic generating function for integers supported on the \(i\)th
group, with \(t\) recording their number of distinct prime factors.
Differentiating its finite product and
evaluating at \(t=1/2\) inserts the factor
\(\omega_i(r)(1/2)^{\omega_i(r)-1}\).
The groups are disjoint, so multiplying these marked sums and
dividing by \(\prod_iV_i\) proves \eqref{con:harmonic-identity}.

Integers with one prime, to exponent one, from each group contribute
\[
 \prod_{i=1}^K\left(V_i^{-1}\sum_{p\in\mathcal P_i}p^{-1}\right)=1
\]
to \(J_0\). For the upper bound and tail put \(s=L^{-a_K}\).
Every group prime satisfies \(p^s\le\e^2\), so, for large \(x\),
\[
 \sum_{j\ge1}p^{-j(1-s)}\le\frac{2\e^2}{p}.
\]
Bounding the mark by \(B_K\) and multiplying the geometric series
over group primes gives
\[
 \sum_r\frac{\mathcal W(r)}{r^{1-s}}\ll_K1,
\]
since \(\sum_iV_i\ll K\).  Multiplication by
\(x^{-\eps s}\) on \(r>x^\eps\) proves the tail.  Finally,
\(A_r\ll x/r\) by counting integers in the interval for \(Q\)
whenever \(A_r\ne0\); otherwise the bound is immediate.
\end{proof}

Take temporarily
\begin{equation}\label{con:small-parameters}
 0<\kappa<0.01,\qquad 0<\eps<\kappa,
 \qquad 0<b<0.01,
 \qquad D=x^{1/2-\kappa/2}.
\end{equation}
For an odd squarefree integer \(\ell\), let \(A_r(\ell)\) denote
the sum defining \(A_r\) with the additional condition
\(\ell\mid crQ+1\), and set
\begin{equation}\label{con:local-density}
 g_r(\ell)=\frac{\one_{\{(\ell,Mr)=1\}}}{\varphi(\ell)},\qquad
 E_r(\ell)=A_r(\ell)-g_r(\ell)A_r.
\end{equation}
In particular, \(g_r(1)=1\) and \(E_r(1)=0\).

\begin{lemma}[Distribution of the weighted family]\label{con:distribution}
For every fixed \(A>0\),
\begin{equation}\label{con:bv-sum}
 \sum_{r\le x^\eps}\mathcal W(r)
 \sum_{\substack{\ell\le D\\\ell\text{ odd squarefree}}}|E_r(\ell)|
       \ll xJ_0L^{-A}.
\end{equation}
Moreover, writing \(A_\Psi=\int_1^2\Psi(y)\,dy\), we have
\begin{equation}\label{con:mass-asymptotic}
 X_0\sim\frac{xJ_0A_\Psi}{c\varphi(M/c)L}
 \qquad(x\to\infty),
\end{equation}
with all construction parameters fixed and \(J_0=J_0(x)\).
In particular,
\begin{equation}\label{con:total-mass}
 X_0\asymp_{M,\Psi}\frac{xJ_0}{L},
\end{equation}
with comparison constants independent of \(K,\kappa,\eps,b\).
\end{lemma}

\begin{proof}
Put \(N=M/c\).  For \(r\le x^\eps\), the support of \(\Psi\)
makes \(Q>x^{0.9}\) automatic.  The progression for \(d\) is
equivalent to the reduced class
\[
 Q\equiv\frac{u-1}{c}\,r^{-1}\pmod N.
\]
If \((\ell,Mr)>1\), then \(A_r(\ell)=g_r(\ell)=0\).
Otherwise the additional condition is
\(Q\equiv-(cr)^{-1}\pmod\ell\), and the two classes combine to
one reduced class modulo \(N\ell\).

We use the prime-counting form of the Bombieri--Vinogradov theorem
\cite{Bombieri65,Vinogradov65}: for fixed \(A',\eta>0\),
\begin{equation}\label{con:bv-input}
 \sum_{q<Y^{1/2-\eta}}\max_{(v,q)=1}
 \left|\pi(Y;q,v)-\frac{\Li(Y)}{\varphi(q)}\right|
       \ll_{A',\eta}Y(\log Y)^{-A'}.
\end{equation}
Here \(\Li(Y)=\int_2^Ydt/\log t\).  Uniformly for
\(r\le x^\eps\) and \(Y\asymp x/(cr)\), the moduli
\(N\ell\), \(\ell\le D\), lie in this range.  Indeed, with
\(\eta=(\kappa-\eps)/4\),
\[
 (1-\eps)(1/2-\eta)-(1/2-\kappa/2)
       =\frac{(\kappa-\eps)(1+\eps)}4>0,
\]
which also absorbs the fixed factor \(N\).

Define
\[
 I_r=\frac{x}{cr}\int_1^2
               \frac{\Psi(y)}{\log((xy-1)/(cr))}\,dy.
\]
For every fixed \(0<\theta<0.1\), the support of \(\Psi\) forces
\(Q>x^{0.9}\) whenever \(r\le x^\theta\). The prime number theorem
in the fixed reduced progressions modulo \(N\), followed by partial
summation, therefore gives, for every fixed \(A>0\),
\begin{equation}\label{con:fixed-r-mass}
 A_r=\frac{I_r}{\varphi(N)}
      +O_{A,M,\Psi}\bigl((x/r)L^{-A}\bigr)
 \qquad(r\le x^\theta).
\end{equation}
The class varies with \(r\), but \(N\) is fixed and every class is
reduced. The implied constant is independent of \(\theta\), although
the threshold for the automatic cutoff may depend on \(\theta\).
This single-modulus estimate thus holds throughout
\(0<\theta<0.1\); the restriction \(\eps<\kappa\) is needed for
the joint modulus estimate that follows.

Partial summation in \eqref{con:bv-input} yields
\begin{equation}\label{con:smoothed-bv}
 \sum_{\substack{\ell\le D\\\ell\text{ odd squarefree}}}
 \left|A_r(\ell)-\frac{g_r(\ell)}{\varphi(N)}I_r\right|
       \ll_{A'}\frac xrL^{-A'}.
\end{equation}
To check the smoothing uniformly, integrate the progression remainder
against the derivative of \(\Psi((crt+1)/x)\).  Its total absolute
integral is \(\int|\Psi'(y)|\,dy\), independently of \(r\), and
all integration points have \(t\asymp x/(cr)\).  The moduli
\(N\ell\) are distinct, and the maximum in \eqref{con:bv-input}
covers their varying residue classes.

The term \(\ell=1\) of \eqref{con:smoothed-bv} evaluates \(A_r\).
Replacing \(I_r/\varphi(N)\) by \(A_r\) in the other terms costs
at most one logarithmic factor, because
\[
 \sum_{\ell\le D}g_r(\ell)\le\sum_{\ell\le D}\frac1{\varphi(\ell)}
       \ll L.
\]
For example, expand
\(n/\varphi(n)=\sum_{d\mid n}\mu^2(d)/\varphi(d)\) and use
\(\sum_d\mu^2(d)/(d\varphi(d))<\infty\).
Choosing \(A'\) sufficiently large gives
\(\sum_{\ell\le D}|E_r(\ell)|\ll(x/r)L^{-A}\); summing with
the marks proves \eqref{con:bv-sum}.

On the support of \(\Psi\), for \(r\le x^\eps\) the logarithm
in \(I_r\) is comparable to \(L\), with absolute comparison
constants since \(\eps<0.01\) and \(c\in\{2,4\}\).  Thus
\(I_r\asymp_\Psi x/(rL)\).  Summing the evaluation of \(A_r\)
and using Lemma~\ref{con:harmonic} proves \eqref{con:total-mass}.
The tail divided by \(xJ_0/L\) tends to zero since \(J_0\ge1\).
Although its threshold may depend on \(K\), the resulting comparison
constants do not.

To obtain the precise asymptotic, fix an auxiliary exponent
\(0<\theta<\eps\), without changing the sieve parameters.
Uniformly for \(r\le x^\theta\) and \(y\) in the support of \(\Psi\),
\[
 \log((xy-1)/(cr))=L+O(\theta L+1),
\]
and hence
\[
 I_r=\frac{xA_\Psi}{crL}
             \bigl(1+O(\theta)+O(L^{-1})\bigr).
\]
The constants in these relative errors are independent of \(\theta\).
Use the \(\ell=1\) term of \eqref{con:smoothed-bv} with \(A'>2\),
multiply by \(\mathcal W(r)\), and sum over \(r\le x^\theta\).
The summed progression error is \(O(xJ_0L^{-A'})\).
Lemma~\ref{con:harmonic} bounds the omitted parts of both \(X_0\)
and \(J_0\); after the harmonic tail is multiplied by \(x/L\),
both are \(o_\theta(xJ_0/L)\), since \(J_0\ge1\). Thus
\[
 \frac{c\varphi(N)L X_0}{xJ_0A_\Psi}
       =1+O(\theta)+o_\theta(1).
\]
First let \(x\) tend to infinity with \(\theta\) fixed, and then
let \(\theta\) decrease to zero. This proves
\eqref{con:mass-asymptotic}. For \((M,c,u)=(8,4,5)\), it becomes
\(X_0\sim xJ_0A_\Psi/(4L)\), since \(\varphi(M/c)=\varphi(2)=1\).
\end{proof}

To express the sieve products, define
\begin{equation}\label{con:singular-product}
 \begin{split}
 V_M(y)&=\prod_{\substack{2<p\le y\\p\nmid M}}
                            \left(1-\frac1{p-1}\right),\\
 \mathfrak S_M&=2\prod_{p>2}\left(1-\frac1{(p-1)^2}\right)
       \prod_{\substack{p\mid M\\p>2}}
                            \left(1-\frac1{p-1}\right)^{-1}>0.
 \end{split}
\end{equation}
Division by \(V(y)\) and Mertens' theorem give
\begin{equation}\label{con:mertens}
 V_M(y)\sim\mathfrak S_M V(y),\qquad
 V(y)\sim\frac{\e^{-\gamma_E}}{\log y},
\end{equation}
where \(\gamma_E\) is Euler's constant.  The primes dividing any
supported \(r\) are large enough that their omission changes these
products negligibly.  More precisely, \eqref{con:support} gives
\begin{equation}\label{con:omitted-primes}
 \sum_{p\mid r}\frac1p\le0.11L^{0.9}\exp(-L^{0.1})=o(1),
 \qquad
 \prod_{2<p\le y}(1-g_r(p))=(1+o(1))V_M(y),
\end{equation}
uniformly for \(r\le x^{0.11}\) and all \(y\).  The bound on the
relative error is independent of \(K\).

\subsection{Two sieve estimates}

The block sieve of Lemma~\ref{lem:block}, from
\cite[Lemma~2.9]{BlockPaper}, now gives the initial mass without
small prime factors and a conditional mass for a specified divisor.
Both estimates use the distribution already proved, but at different
sieve depths.

First sieve the prime objects \(Q\) in \(A_r\), for
\(r\le x^\eps\), by the conditions \(p\mid crQ+1\) at odd
primes \(p\le z=x^b\).  The local densities are \(g_r(p)\),
with \(g_r(p)\le1/2\) and
\(\sum_{v<p\le v^2}g_r(p)\ll1\) with absolute constants.
Choose the even depth
\[
 H=2\left\lfloor\frac1{40b}\right\rfloor.
\]
For sufficiently small fixed \(b\), this is large enough for the
lower bound in Lemma~\ref{lem:block}; moreover,
\[
 z^{4H+2}\le x^{0.2+2b}<D,\qquad
 \e^{-H}\le\e^2\exp(-1/(20b)).
\]
All remainders are therefore covered by \eqref{con:bv-sum}.
Every supported \(d\) is odd.  Summing the sieve bound with the
marks, using \eqref{con:mertens}--\eqref{con:omitted-primes} and
the negligible tail in Lemma~\ref{con:harmonic}, gives
\begin{equation}\label{con:initial-lower}
 S_b:=\sum_{P^-(d)>x^b}w(d)
 \ge\frac{\mathfrak S_M X_0}{L}
 \left\{\frac{\e^{-\gamma_E}}b
              (1-C_s\e^{-c_s/b})+o(1)\right\},
\end{equation}
where \(C_s,c_s>0\) are absolute; the displayed bound on \(\e^{-H}\)
permits \(c_s=1/20\), after enlarging \(C_s\).
The remainders are
\(o(X_0/L)\) by taking \(A>2\) in \eqref{con:bv-sum}.

\begin{lemma}[Mass conditioned on a divisor]\label{con:bin-mass}
For any fixed \(b\le\gamma<\gamma'\le1/2-\kappa\),
\begin{equation}\label{con:bin-mass-formula}
 \sum_{\substack{x^\gamma<n\le x^{\gamma'}\\n\text{ prime}}}
 \sum_{\substack{n\mid d\\P^-(d)>W}}w(d)
 =\mathfrak S_M X_0 V(W)
                  \left(\log\frac{\gamma'}\gamma+o(1)\right).
\end{equation}
\end{lemma}

\begin{proof}
For large \(x\), the primes \(n\) in this sum exceed \(W\) and
all group primes, and do not divide \(M\).  Fix \(r\le x^\eps\)
and condition the family in \(A_r\) on \(n\mid crQ+1\).  For a
squarefree product \(\ell\) of odd primes at most \(W\), the
intersection count is
\[
 A_r(n\ell)=\frac{A_r}{n-1}g_r(\ell)+E_r(n\ell).
\]
This includes \(\ell=1\), as allowed in Lemma~\ref{lem:block}.
Apply that lemma with main mass \(A_r/(n-1)\) and
\(H=2\lceil(\log L)^2\rceil\).  Its relative error tends to zero,
and its remainder moduli satisfy
\[
 \ell\le W^{4H+2}=\exp(O(L^{0.24}(\log L)^2))=x^{o(1)},
 \qquad n\ell\le x^{1/2-\kappa+o(1)}<D.
\]
The map \((n,\ell)\mapsto n\ell\) is injective, since \(n\) is
the unique prime factor exceeding \(W\).  Thus the total remainder
after summing \(n\) and the marked \(r\)'s is bounded by a single
copy of \eqref{con:bv-sum}.

The sieve product is \((1+o(1))V_M(W)\) uniformly in \(r\), and
Mertens' theorem gives
\[
 \sum_{\substack{x^\gamma<n\le x^{\gamma'}\\n\text{ prime}}}
                      \frac1{n-1}=\log(\gamma'/\gamma)+o(1).
\]
By \eqref{con:mass-definitions}, Lemma~\ref{con:harmonic}, and
\eqref{con:total-mass}, for every fixed \(A>0\) we have
\[
 \sum_{r\le x^\eps}\mathcal W(r)A_r
       =X_0+o(X_0L^{-A}).
\]
This evaluates the contribution of \(r\le x^\eps\) with the exact
main mass \(X_0\); taking \(A\) large in \eqref{con:bv-sum}
makes the remainders \(o(X_0V(W))\).  Finally, each supported \(d\)
has at most \(\log(2x)/(bL)=O(1/b)\) distinct prime divisors above
\(x^b\).  Consequently the harmonic tail bounds the omitted
\(r>x^\eps\) contribution even after summing over \(n\).
This proves \eqref{con:bin-mass-formula}.
\end{proof}

\subsection{Removing composites by their least prime factor}

We next pass from \(S_b\) to
\[
 U_\kappa=\sum_{P^-(d)>x^{1/2-\kappa}}w(d).
\]
For a bin \((\gamma,\gamma']\subseteq(b,1/2-\kappa]\), with the
first left endpoint allowed to equal \(b\), define
\[
 T_{\gamma,\gamma'}=
 \sum_{\substack{x^\gamma<n\le x^{\gamma'}\\n\text{ prime}}}
                    \sum_{P^-(m)>x^\gamma}w(mn).
\]
If \(x^\gamma<P^-(d)\le x^{\gamma'}\), taking \(n=P^-(d)\)
and \(m=d/n\) shows that its weight is included in this sum.
This remains true when \(n^2\mid d\).

Partition both factors into disjoint half-open dyads with endpoints
powers of two, retaining pairs that meet the smooth product support
and the range of \(n\).  On that support, with \(t=\log n/L\),
\[
 \frac{\log m}{L}=1-t+O(L^{-1}).
\]
Every retained full \(m\)-dyad therefore satisfies
\eqref{con:exponent-range} with
\[
 s_-=1/2+\kappa/2,\qquad s_+=1-b/2,
\]
for large \(x\); in particular \(s_->\gamma\).
Both dyad scales are at least \(x^{b/2}\), and their product is
comparable to \(x\).  There are \(O(L)\) such pairs, since each
\(n\)-dyad allows only a bounded number of \(m\)-dyads.

Apply Lemma~\ref{con:type-ii} in every pair with \(\delta=b/2\)
and with \(b_n\) the prime indicator restricted to its dyad and
bin.  These coefficients are \(W\)-rough.  The total error is
\(O(xL^{-5})=o(X_0/L)\), by \eqref{con:total-mass} and
\(J_0\ge1\).  Hence
\[
 T_{\gamma,\gamma'}=
 \sum_{\substack{x^\gamma<n\le x^{\gamma'}\\n\text{ prime}}}
 \sum_{m>x^\gamma}\frac{D_\gamma(\log m/L)}{LV(W)}
                         \one_{\{P^-(m)>W\}}w(mn)+o(X_0/L).
\]
Joint continuity of the density on the fixed compact exponent range
gives, on the support of the summand,
\[
 D_\gamma(\log m/L)
       \le\sup_{\gamma\le t\le\gamma'}D_\gamma(1-t)+o(1).
\]
Since \(n>W\), the remaining roughness condition is precisely
\(P^-(mn)>W\).  Lemma~\ref{con:bin-mass} thus gives
\begin{equation}\label{con:bin-cost}
 T_{\gamma,\gamma'}\le\frac{\mathfrak S_M X_0}{L}
 \left\{\sup_{\gamma\le t\le\gamma'}D_\gamma(1-t)
                       \log\frac{\gamma'}\gamma+o(1)\right\}.
\end{equation}

The density estimate in Lemma~\ref{lem:density} is normalized to
match this expression:
\begin{equation}\label{con:density-budget}
 \int_b^{1/2}D_t(1-t)\frac{dt}{t}=D_b(1)-1,
 \qquad
 D_b(1)\le\frac{\e^{-\gamma_E}}b(1+C_d\e^{-c_d/b})
\end{equation}
for sufficiently small fixed \(b\). On
\(b\le\gamma\le t\le1/2-\kappa\), the point \((\gamma,1-t)\)
stays a fixed distance from \(s=\gamma\).  Uniform continuity and
the finite mass of \(dt/t\) therefore give a sufficiently fine
fixed partition
\[
 b=\gamma_0<\gamma_1<\cdots<\gamma_J=1/2-\kappa
\]
such that
\begin{align*}
 \sum_{j=1}^J
 \sup_{\gamma_{j-1}\le t\le\gamma_j}D_{\gamma_{j-1}}(1-t)
                         \log\frac{\gamma_j}{\gamma_{j-1}}
 &\le\int_b^{1/2-\kappa}D_t(1-t)\frac{dt}{t}+0.1\\
 &\le D_b(1)-1+0.1.
\end{align*}
Subtracting \eqref{con:bin-cost} for these finitely many bins from
\eqref{con:initial-lower} proves
\begin{equation}\label{con:post-bin}
 U_\kappa\ge\frac{\mathfrak S_M X_0}{L}
 \left\{0.9-\frac{\e^{-\gamma_E}}b
       (C_s\e^{-c_s/b}+C_d\e^{-c_d/b})+o(1)\right\}.
\end{equation}
The unit term in \eqref{con:density-budget} is what leaves a positive
mass after this subtraction.  It remains to remove the composites
whose two prime factors are both close to \(\sqrt x\).

\subsection{Composites with two nearly equal prime factors}

A composite counted by \(U_\kappa\) has exactly two prime factors
with multiplicity: three factors greater than \(x^{1/2-\kappa}\)
would have product exceeding \(2x\).  Each factor lies in
\([x^{1/2-\kappa},2x^{1/2+\kappa}]\).  Cover the ordered factor
pairs by dyadic boxes
\begin{equation}\label{con:boxes}
 \mathcal B=[M_1,2M_1)\times[M_2,2M_2)
\end{equation}
that meet this range in each coordinate and admit a product in
\([x,2x]\).  They satisfy \(x/4\le M_1M_2\le2x\), and their
full coordinate intervals lie in \([x^{0.46},x^{0.54}]\) for large
\(x\).  There are at most \(C_{\rm dyad}(\kappa L+1)\) boxes,
where \(C_{\rm dyad}\) is absolute: this bounds the number of first
coordinate dyads, and the product condition leaves only a bounded
number of second coordinate dyads for each.

The rough-pair bound holds on a slightly wider fixed range of boxes,
which includes all the boxes just selected.
\begin{lemma}[An upper bound for rough pairs]\label{con:rough-pairs}
Let \(\mathcal B=[M_1,2M_1)\times[M_2,2M_2)\) be any dyadic box
such that
\begin{equation}\label{con:rough-box-range}
 x/4\le M_1M_2\le2x,\qquad
 x^{0.46}/2\le M_i\le2x^{0.54}\quad(i=1,2).
\end{equation}
Then
\begin{equation}\label{con:rough-pair-bound}
 \sum_{\substack{(m,n)\in\mathcal B\\P^-(m)>W,\ P^-(n)>W}}w(mn)
       \le C_1X_0V(W)^2,
\end{equation}
where \(C_1\) depends at most on \(M,\Psi\), and is independent
of \(K,\kappa,b,\eps\).
\end{lemma}

\begin{proof}
Fix a contributing group integer \(r\le x^{0.11}\), and set
\(k_r=cr\), \(z'=x^{0.002}\).  The corresponding pairs satisfy
\(mn\equiv1\pmod{k_r}\).  At every odd prime \(j\le z'\) with
\(j\nmid r\), they avoid
\begin{equation}\label{con:pair-bad}
 mn\equiv1\pmod j,
 \qquad\text{and, if }j\le W,\qquad mn\equiv0\pmod j.
\end{equation}
Indeed, \(mn-1=k_rQ\) with \(Q>x^{0.9}>z'\), while both factors
are \(W\)-rough.  We impose no further condition at primes dividing
\(k_r\) and may omit the fixed progression for this upper bound.

For a squarefree product \(\ell\) of the designated odd primes,
put
\[
 \rho(\ell)=\prod_{j\mid\ell}
               \bigl(j-1+(2j-1)\one_{\{j\le W\}}\bigr).
\]
The conditions \(mn=1\) and \(mn=0\) modulo \(j\) are disjoint
and have respectively \(j-1\) and \(2j-1\) residue pairs.  There
are \(\varphi(k_r)\) pairs with product one modulo \(k_r\).
The Chinese remainder theorem therefore gives exactly
\(\varphi(k_r)\rho(\ell)\) pairs modulo \(k_r\ell\) satisfying
the base congruence and all the forbidden conditions at primes
dividing \(\ell\).  Their count in the box is
\[
 X'g(\ell)+E(\ell),\qquad
 X'=M_1M_2\frac{\varphi(k_r)}{k_r^2},\qquad
 g(\ell)=\frac{\rho(\ell)}{\ell^2},
\]
where
\begin{equation}\label{con:lattice-error}
 E(\ell)\ll\varphi(k_r)\rho(\ell)
             \left(\frac{M_1+M_2}{k_r\ell}+1\right).
\end{equation}
This formula also holds at \(\ell=1\).

The local bounds \(g(j)\le7/9\) and \(g(j)\le3/j\) verify
the gap and block-sum hypotheses of Lemma~\ref{lem:block}, with
absolute constants.  Its upper bound with \(H=2\) uses
\(\ell\le(z')^{10}=x^{0.02}\). For squarefree \(\ell\),
\[
 \rho(\ell)\le\prod_{j\mid\ell}3j\le\tau(\ell)^2\ell,
\]
where \(\tau\) is the divisor function. Since \(\varphi(k_r)\le k_r\)
and \(k_r=cr\le4r\), \eqref{con:lattice-error} gives
\[
 |E(\ell)|\ll\tau(\ell)^2(M_1+M_2+r\ell).
\]
The needed divisor sum follows elementarily. Let \(\tau_4(n)\)
count ordered decompositions of \(n\) into four positive factors.
Prime by prime, \(\tau(n)^2\le\tau_4(n)\); summing over the
first three factors therefore yields, for \(Y\ge2\),
\[
 \sum_{n\le Y}\tau(n)^2
 \le Y\left(\sum_{a\le Y}\frac1a\right)^3
 \ll Y\bigl(\log(2Y)\bigr)^3.
\]
Consequently, taking \(Y=x^{0.02}\) and using
\(M_1+M_2\le4x^{0.54}\) from \eqref{con:rough-box-range}, we obtain
\begin{equation}\label{con:lattice-sum}
 \sum_{\ell\le x^{0.02}}|E(\ell)|
       \ll (x^{0.56}+rx^{0.04})L^3
       =o\bigl((x/r)L^{-10}\bigr)
\end{equation}
uniformly for \(r\le x^{0.11}\).  The ratios of the two displayed
power terms to \(x/r\) are at most \(x^{-0.33}\) and
\(x^{-0.74}\), respectively, and absorb the logarithmic factor.

For \(j\le W\), the surviving sieve factor is
\[
 1-\frac{3j-2}{j^2}
       =\left(1-\frac1j\right)^3
                         \left(1-\frac1{(j-1)^2}\right),
\]
and for \(W<j\le z'\) it is
\[
 1-\frac{j-1}{j^2}
       =\left(1-\frac1j\right)
                         \left(1+\frac1{j(j-1)}\right).
\]
The correction products are bounded.  Omitting the primes dividing
\(r\) costs a bounded factor by \eqref{con:omitted-primes}, and
omitting two costs a fixed factor.  Thus
\[
 \prod_{\substack{2<j\le z'\\j\nmid r}}(1-g(j))
       \ll V(W)^2V(z').
\]
As \(X'\ll x/r\), the upper sieve and \eqref{con:lattice-sum}
bound the pair count at this \(r\) by
\(O((x/r)V(W)^2V(z'))\), uniformly in all the marks.  Bound
\(\Psi\) by one, multiply by \(\mathcal W(r)\), and sum.  The
result is
\[
 O\bigl(xJ_0V(W)^2V(z')\bigr).
\]
Finally \(V(z')\asymp L^{-1}\), with fixed comparison constants,
and \eqref{con:total-mass} proves \eqref{con:rough-pair-bound}.
Only \(J_0\), rather than the pointwise bound \(B_K\), entered this
calculation, so \(C_1\) has the claimed independence.
\end{proof}

Return now to the boxes selected in \eqref{con:boxes}, whose full
coordinate intervals lie in \([x^{0.46},x^{0.54}]\). We replace
the prime indicators in each of these boxes. Apply
Lemma~\ref{con:type-ii} with
\[
 \gamma=0.4,\qquad s_-=0.46,\qquad s_+=0.54,
 \qquad\delta=b/2.
\]
In this range \(R_{0.4}(m)\) is exactly the prime indicator, since
\(2(0.4)>0.54\), and \(D_{0.4}(s)=1/s\).  Its comparison
coefficient is therefore
\[
 B_{\rm pr}(m)=\frac{\one_{\{P^-(m)>W\}}}{V(W)\log m}.
\]
The identity
\[
 \begin{aligned}
 \one_{\{m\text{ prime}\}}\one_{\{n\text{ prime}\}}
       -B_{\rm pr}(m)B_{\rm pr}(n)
 &=(\one_{\{m\text{ prime}\}}-B_{\rm pr}(m))\one_{\{n\text{ prime}\}}\\
 &\quad+B_{\rm pr}(m)(\one_{\{n\text{ prime}\}}-B_{\rm pr}(n))
 \end{aligned}
\]
allows two successive applications, interchanging the variables in
the second.  The companion coefficient is \(W\)-rough and bounded
by one for large \(x\) in both cases.  All coefficients retain
their box restrictions.  The total error over the boxes is
\(O(xL^{-5})=o(X_0/L)\).

Since \(\log m,\log n\ge0.46L\), Lemma~\ref{con:rough-pairs}
bounds the resulting sum in one box by
\((0.46)^{-2}C_1X_0/L^2\).  The ordered prime-pair sum includes
prime squares and may count other semiprimes twice, which is harmless
for an upper bound.  Summing the boxes gives
\begin{equation}\label{con:balanced-cost}
 \sum_{\substack{P^-(d)>x^{1/2-\kappa}\\d\text{ composite}}}w(d)
       \le C_2(\kappa+L^{-1})\frac{X_0}{L}+o(X_0/L),
\end{equation}
where \(C_2=(0.46)^{-2}C_1C_{\rm dyad}\) depends at most on
\(M,\Psi\), and is independent of \(K,\kappa,b,\eps\).

\subsection{Choice of parameters and completion}

\begin{proof}[Completion of the proof of Proposition~\ref{prop:controlled}]
Fix \(M,c,u,\Psi\) first, and hence the constants
\(\mathfrak S_M\) and \(C_2\).  Choose
\(0<\kappa<0.01\) such that
\[
 C_2\kappa<\mathfrak S_M/4,
\]
and then choose \(0<\eps<\kappa\).  Choose \(b>0\) sufficiently
small for the sieve and density estimates, and so that
\[
 \frac{\e^{-\gamma_E}}b
          (C_s\e^{-c_s/b}+C_d\e^{-c_d/b})<0.2.
\]
Fix the partition used in the least-factor subtraction.  Only
finitely many Type~II parameter sets now occur.  Choose \(K\)
sufficiently large for all of them, and choose the exponents
afterward, for example
\[
 a_i=0.1+\frac{0.1i}{K+1}\qquad(1\le i\le K).
\]
This order is permitted by Lemma~\ref{lem:marked-coefficient-transfer}
and Proposition~\ref{prop:rough-type-ii}, hence by Lemma~\ref{con:type-ii}.

Every supported prime is counted by \(U_\kappa\).
Subtracting \eqref{con:balanced-cost} from \eqref{con:post-bin},
and then taking \(x\) sufficiently large for the fixed parameters,
gives
\[
 \frac{L}{\mathfrak S_M X_0}\sum_{d\text{ prime}}w(d)
       \ge0.9-0.2-\frac14+o(1)>\frac14.
\]
By \eqref{con:total-mass}, \(J_0\ge1\), and \(w\le B_K\),
the number of distinct supported primes is therefore
\(\gg_{M,c,u}x/L^2\).  Their support in \eqref{con:weight}
ensures \(x<d<2x\), \(d\equiv u\pmod M\), and
\(d=crQ+1\) with \(Q>x^{0.9}\) prime.  Every prime factor of
\(r\) belongs to a group in \eqref{eq:marked-prime-groups}, hence lies in
\((\exp(L^{0.1}),\exp(L^{0.3}))\) for all sufficiently large
\(x\).  These are precisely the assertions of the proposition.
\end{proof}
